\documentclass[a4paper,10pt]{report}

\def\packagepath{../../preambule}   % path du package princial
\usepackage{\packagepath/preambule}    % utilisation du fichier de configuration

\def\level{Première }              % Classe
\def\course{NSI}              % Matière
\def\eval{}

\def\visibleornot{visible}    % visible or invisible
\def\documentpath{./Sources_Latex}
\def\date{Mardi 16 Décembre 2025}
\renewcommand{\arraystretch}{1}  % Ecart dans les tableau

\def\ptsexoA{4}
\def\ptsexoB{2}
\def\ptsexoC{3}
\def\ptsexoD{1}
\def\ptsexoE{2}
\def\ptsexoF{4}
\def\ptsexoG{4}
\def\ptsexoH{2}
\def\ptsexoI{2}
\def\ptsexoJ{2}
\def\ptsexoK{2}

%\def\ptstotal{\ptsexoA+\ptsexoB+\ptsexoC+\ptsexoD}
%\usepackage{pgf}
\pgfmathtruncatemacro{\ptstotal}{\ptsexoA+\ptsexoB+\ptsexoC+\ptsexoD}
\usepackage{hyperref}
\usepackage{tasks}
\usetikzlibrary{shapes, arrows, positioning}
\usetikzlibrary{shapes.geometric, arrows, positioning, decorations.pathreplacing}

\tikzstyle{etat} = [draw, rounded corners=15pt, minimum width=2.5cm, minimum height=1.2cm, text centered, font=\bfseries]
\tikzstyle{nouveau} = [etat, fill=blue!40]
\tikzstyle{pret} = [etat, fill=green!60]
\tikzstyle{elu} = [etat, fill=yellow!60]
\tikzstyle{bloque} = [etat, fill=red!60]
\tikzstyle{termine} = [etat, fill=white]
\tikzstyle{fleche} = [thick, ->, >=stealth]

\usepackage{float} % Permet d'utiliser [H]
\usepackage[T1]{fontenc}
%\newcommand{\field}[2]{\TextField[width=#2]{#1}}

%\newcommand{\Champ}[3][]{%
    %\TextField[name=#2,width=#3,height=1.2em,#1]{}
%}

\begin{document}


\begin{Form}

\begin{center}
%{\Large\bfseries Devoir Surveillé de Mathématiques}\\[0.5cm]
{\large 
Nom : \TextField[name=nom, width=4cm, bordercolor=red, backgroundcolor={0.85 0.90 1}]{} \hfill 
Prénom : \TextField[name=prenom, width=4cm, bordercolor=red, backgroundcolor={0.85 0.90 1}]{} \hfill 
Note : \TextField[name=notetotale, width=2cm, bordercolor={0 0 0}, backgroundcolor={0.9 0.9 0.9}, readonly=true, calculate={%
var sum = 0;
var count = 0;
for(var i = 1; i <= 10; i++) {
  var field = this.getField("note" + i);
  if(field && field.value !== "") {
    var val = parseFloat(field.value);
    if(!isNaN(val)) {
      sum += val;
      count++;
    }
  }
}
if(count > 0) {
  event.value = sum + " / \ptstotal";
} else {
  event.value = "";
}
}]{} }
\end{center}


%%%%%%%%%%%%%%%%%%%%%%%%%%%%%%%%%%%%%%%%%%%%%%%%%%ù

%\PageGardeSujetBac[Matiere = Numérique et Sciences Informatiques,
%                  Session = 2025,
%                  AffJour=false,
%                  Duree = {3 heures 30},
%                  DernierePage = \pageref{LastPage},
%                  ModeExamen=false
%                  ]
\renewcommand{\labelitemi}{\textbullet} %pour éviter les tirets dans les "itemize" qui apportent confusion avec le signe moins.
%% Début page de garde style BAC

   %%%%%%%%%%%%%%%%%%%%%%%%%%%%%%%%%%%%%%%%%%%%%%%%%%%%%%%%
%\PageGardeSujetBac[clés]
\pagestyle{DS_FP}          % Feuille de style fancy
%\NomPrenomNote{}

%\cadreNom{}
\appreciation{}

\begin{center}
\textbf{ ------ Consignes ------}
\end{center}

\medskip

\begin{enumerate}[a)]
    \item Vérifier que vous avez bien à l'intérieur les trois fichiers suivants:
    \begin{itemize}
        \item \verb|1NSI_internet.pdf|
        \item \verb|1NSI_reseau01.pkt|
        \item \verb|1NSI_reseau02.pkt|
    \end{itemize}
    \item Dans la barre de recherche de Windows chercher le logiciel nommé \verb|PDFXViewer|
    \item Ouvrir ce logiciel.
    \item \verb|Fichier/Ouvrir| et ouvrir le fichier \verb|1NSI_internet.pdf|. 
    
    (Ce PC $\rightarrow$ \verb|P:\ Espace personnel de xxx \ 1NSI-DSinternet| )
    \item Vous pouvez alors modifier ce fichier PDF. N'oubliez pas de sauvegarder régulièrement.
\end{enumerate}


\bigskip

%Exercice 1
\exodst{\ptsexoA}



Donner une définition des termes suivants:

\begin{enumerate}
    \item Adresse IP (ou adresse Réseau) \hfill \Bareme{} 
    
\longfield{3cm}

\item Masque de sous réseau \hfill \Bareme{} 
    
\longfield{3cm}

\item Serveur DHCP (Dynamic Host Configuration Protocol) \hfill \Bareme{}

\longfield{3cm}

\item Serveur DNS (Domain Name Systeme) \hfill \Bareme{}

\longfield{3cm}
\end{enumerate}


\medskip

%Exercice 2
\exodst{\ptsexoB}

Indiquer par une croix dans la case correspondante si les adresses IP suivantes peuvent être données à des appareils connectés::
    
\renewcommand{\arraystretch}{1.5}
\begin{tabular}{|p{6cm}|c|c|}
    \hline
    \textbf{Adresse IP} & \textbf{Peut être donnée} & \textbf{Ne peut pas être donnée} \\
    \hline
    \verb|192.168.1.4.2| & \checkbox{} & \checkbox{} \\
    \hline
    \verb|192.168.1.4| & \checkbox{} & \checkbox{} \\
    \hline
    \verb|10.255.255.254| & \checkbox{} & \checkbox{} \\
    \hline
    \verb|192.168.1.255| & \checkbox{} & \checkbox{} \\
    \hline
\end{tabular}

\medskip


%Exercice 3
\exodst{\ptsexoC}

Voici la configuration réseau d'un appareil:

\begin{center}
    \includegraphics[width = 0.6\linewidth]{ipconfig2.png}
\end{center}

\begin{enumerate}
    \item Indiquer l'adresse IP de cet appareil. \smallfield{} \hfill \Bareme{}
    \item Indiquer le masque de sous-réseau de cet appareil. \smallfield{} \hfill \Bareme{}
    \item Indiquer l'adresse IP de la passerelle par défaut. \smallfield{} \hfill \Bareme{}
    \item Indiquer l'adresse MAC de cet appareil. \smallfield{} \hfill \Bareme{}
    \item Indiquer l'adresse de ce réseau. \smallfield{} \hfill \Bareme{}
    \item Indiquer l'adresse de Broadcast de ce réseau. \smallfield{} \hfill \Bareme{}
\end{enumerate}

\medskip

%Exercice 4
\exodst{\ptsexoD}

On suppose que nous avons le réseau suivant:
\begin{itemize}
    \item adresse IP de ce réseau: \verb|192.168.1.0| avec un masque de sous réseau \verb|255.255.255.0|
    
\end{itemize}

\begin{enumerate}
    \item Indiquer la première adresse IP attribuable à un appareil de ce réseau. \smallfield{} \hfill \Bareme{}
    \item Indiquer la dernière adresse IP attribuable à un appareil de ce réseau. \smallfield{} \hfill \Bareme{}
    \item Indiquer le nombre d'appareil que l'on peut intégrer dans ce réseau. \smallfield{} \hfill \Bareme{}
\end{enumerate}

\medskip
\pagebreak

%Exercice 5
\exodst{\ptsexoE}

On suppose que nous avons le réseau suivant:
\begin{itemize}
    \item adresse IP de ce réseau: \verb|172.50.128.0/18|
    
\end{itemize}

\begin{enumerate}
    \item indiquer le masque de sous réseau en notation décimale pointée. (4 valeurs) \smallfield{} \hfill \Bareme{}
    \item Indiquer la première adresse IP attribuable à un appareil de ce réseau. \smallfield{} \hfill \Bareme{}
    \item Indiquer la dernière adresse IP attribuable à un appareil de ce réseau. \smallfield{} \hfill \Bareme{}
    \item Indiquer le nombre d'appareil que l'on peut intégrer dans ce réseau. \smallfield{} \hfill \Bareme{}
\end{enumerate}

\medskip

%Exercice 6
\exodst{\ptsexoF}

Dans cet exercice, nous utiliserons le fichier nommé \verb|1NSI_reseau01.pkt|.

\begin{minipage}{0.6\linewidth}
    \begin{enumerate}[a)]
    \item Ouvrir le logiciel \verb|Packet Tracer|
    \item Ouvrir le fichier \verb|1NSI_reseau01.pkt|
    \item Vous devez obtenir le résultat suivant:
\end{enumerate}
\end{minipage}\hfill
\begin{minipage}{0.37\linewidth}

\includegraphics[width = 0.95\linewidth]{DS03_01.png}
\end{minipage}

\bigskip

\begin{enumerate}
    
    \item Indiquer l'adresse du réseau dans lequel se trouve \verb|PC0|. \smallfield{}\hfill \Bareme{} 
    \item Compléter l'adresse IP et le masque de sous-réseau de \verb|PC2| de manière à faire communiquer \verb|PC0| et \verb|PC2|. 
    
    Inscrire ici l'adresse que vous avez attribuée: \smallfield{} et le masque utilisé: \smallfield{}\hfill \Bareme{}  
    \item Ajouter un routeur modèle \verb|1841| connecté au switch \verb|Switch0|. 
    
\includegraphics[width = 0.2\linewidth]{DS03_pk1.png} \qquad 
\includegraphics[width = 0.2\linewidth]{DS03_pk2.png} \qquad 
\includegraphics[width = 0.3\linewidth]{DS03_pk3.png}

    Vous devriez maintenant obtenir le schéma suivant:

    \begin{center}
        \includegraphics[width = 0.6\linewidth]{DS03_02}
    \end{center} 

    \item Configurer ce routeur de manière à permettre la communication entre \verb|PC0| et le routeur \verb|1841|. On donnera au routeur la dernière adresse IP disponible dans le réseau de \verb|PC0|. (Attention à ne pas oublier d'allumer l'interface réseau du routeur).  
    
    Indiquer ici l'adresse que vous avez attribuée: \smallfield{}  \hfill \Bareme{}  
    \item Vérifier que toutes les machines peuvent communiquer entre elles. \hfill \Bareme{} 
    \item Sauvegarder votre travail.
    
    \bigskip

     
      
        %\appelprof


    


\end{enumerate}

\pagebreak

%Exercice 7
\exodst{\ptsexoG}

\medskip

Dans cet exercice, nous utiliserons le fichier nommé \verb|1NSI_reseau02.pkt|.

    \begin{enumerate}[a)]
    \item Ouvrir le logiciel \verb|Packet Tracer|
    \item Ouvrir le fichier \verb|1NSI_reseau02.pkt|
    \item Vous devez obtenir le résultat suivant:
\end{enumerate}

\includegraphics[width = 0.95\linewidth]{DS03_03.png}


\begin{enumerate}
    \item Vérifier que \verb|PC0| et \verb|PC1| peuvent communiquer entre eux. 
    \item Vérifier que \verb|PC2|, \verb|PC3|, \verb|PC4| et le serveur \verb|DHCP| peuvent communiquer entre eux. 
    \item Les machines du réseau de gauche (\verb|PC0|, \verb|PC1|) ont été configurées pour utiliser une adresse IP fixe (statique) et pour utiliser comme passerelle par défaut le routeur \verb|Routeur0|. Inscrire ici l'adresse IP de la passerelle par défaut en explorant les configurations réseaux des différentes machines du réseau de gauche: \smallfield{}\hfill \Bareme{} 
    \item Les machines du réseau de droite (\verb|PC2|, \verb|PC3|, serveur DHCP) ont été configurées pour utiliser une adresse IP fixe (statique) et pour utiliser comme passerelle par défaut le routeur \verb|Routeur0|. Inscrire ici l'adresse IP de la passerelle par défaut en explorant les configurations réseaux des différentes machines du réseau de droite: \smallfield{}\hfill \Bareme{} 
    \item Configurer les interfaces réseaux du routeur \verb|routeur0| en utilisant les adresses IP trouvées aux questions précédentes.
    \item Vérifier que \verb|PC0| et \verb|PC3| peuvent communiquer entre eux.\hfill \Bareme{} 
    \item Sauvegarder votre travail.
    
    \medskip

    \item Le \verb|PC4| a été configuré en DHCP activé. Indiquer ici son adresse IP: \smallfield{}\hfill \Bareme{} 
    \item Explorer la configuration réseau de cette machine et indiquer comment elle a obtenu son adresse IP.\hfill \Bareme{} 
    
    \longfield{4.5cm}
  

\end{enumerate}

\medskip



\end{Form}

\end{document}










